\section{Discussion}\label{sec:discussion}
% statements-of-record rows resolved here: none

The classical results bearing on the Strong BSD formula each control part
of it.  Gross--Zagier and Kolyvagin~\citep{GrossZagier1986,Kolyvagin1989}
relate derivatives, Heegner points, Selmer groups and regulators in rank
at most one; Iwasawa main conjectures~\citep{SkinnerUrban2014} give
$p$-parts of the formula at many primes; Cassels--Tate pairings and
Selmer complexes~\citep{Flach1990,Nekovar2006} organize the finite
arithmetic.  None of them gives the whole formula in general.  This paper
does not change that.  What it does is make one precise statement about
how the formula fits a closure framework, and record exactly which inputs
such a formulation needs.

The precise statement is the fixed-point formulation of
\Cref{thm:closure:scalar-closure-translation}: the identity holds exactly
when the pair of its two sides is fixed by the averaging closure.  It is
elementary, and it is exact; in particular it separates fixedness of the
original pair from the trivial closedness of its image.
\Cref{thm:closure:formed-closure-bridge} says what a closure on a richer
layer of arithmetic data would have to satisfy to return the identity: a
readout intertwining it with the averaging closure, and fixedness of the
original state.  Neither is constructed here, and idempotence alone does
not supply them.

The main theorem should be read with its hypotheses in view.  Its scalar
conclusion follows from a recognition hypothesis that contains the
identity in rank at most one and expresses it through a normalized
leading-term identity in rank at least two.  The other hypotheses, the
local recognition hypotheses and the comparison inputs, are carried with
their scopes but do not enter the derivation.  The value of the
formulation therefore lies in its accounting: what is assumed, where each
assumption would come from, what its scope is, and what is not known.
Several points of that accounting are not obvious and are worth
recording.  The local congruences modulo $\Zp^\times$ cannot see a
Tamagawa number prime to $p$ (\Cref{sec:recognition-sources}).  The
determinant ideal of a Selmer complex does not determine whether a
prescribed generator has a square root (\Cref{ex:closure:1913b1}).  The
proposed $\eta$-formula involves an eigenvalue of slope $1/2$, not the
critical slope (\Cref{sec:eta-formula}).  And sensitivity of the BSD
product to its factors does not make any hypothesis necessary
(\Cref{rmk:closure:sensitivity-not-necessity}).

The companion paper~\citep{TsiokosBSDApparatus2026} studies the same
factors from the other side: which data each factor of the BSD formula
retains, with finite adequacy tests and examples of information loss.
Its normalization of $\kapr$ is what identifies the rank-$\ge2$
recognition hypothesis with the Strong BSD identity here.

The same pattern, a target statement written as the fixed point of a
closure with the deep inputs named as recognition hypotheses and an
audit register recording their status, is the general method of the Six
Birds framework~\citep{TsiokosFoundationsIII2026,TsiokosAOR2026}.  For BSD
it yields a conditional theorem of exactly the strength of its
hypotheses, together with a precise description of the inputs that an
unconditional proof along these lines would require.
